\section*{Capítulo \ref{ch:b2:reduction} --- Redução de endomorfismos}

\begin{solution}{exo:b2:reduction:1}
$A $: $\chi_A = X^2 - 2X - 3 = (X - 3)(X + 1)$. Autovetores: para
$3$: $(1,1)$; para $-1$: $(1,-1)$. Então $ P = \begin{pmatrix} 1 & 1\\ 1
& -1\end{pmatrix}$ dá $ P^{-1}AP = \operatorname{diag}(3, -1)$.

$ B = J - I $ onde $ J $ é a matriz de todos. $ J $ tem classificação $1$ com
$ Jv = 3v $ para $ v = (1,1,1)$ e $ Jw = 0$ no plano $ x + y + z =
0$: \omterm{def:b2:reduction:eigen}{espectro} de $ B $ é $\{2, -1\}$ com autoespaços
$\operatorname{Vect}(1,1,1)$ (dimensão $1$) e $\{x + y + z = 0\}$
(dimensão $2$, base $(1,-1,0), (1,0,-1)$). $ P $ com esses três
colunas dá $ P^{-1}BP = \operatorname{diag}(2, -1, -1)$.
\end{solution}

\begin{solution}{exo:b2:reduction:2}
\emph{Espaços próprios:} $\chi_C = (X-1)^2$, solteiro \omterm{def:b2:reduction:eigen}{autovalor} $1$;
$\ker(C - I) = \ker\begin{pmatrix} 0&1\\ 0&0\end{pmatrix}$ é o
linha $\operatorname{Vect}(e_1)$: dimensão $1 < 2 = m_1$, então não
\omterm{def:b2:reduction:diag}{diagonalizável} (\cref{thm:b2:reduction:diagcrit}).

\emph{Polinômio mínimo:} $\mu_C $ divide $(X-1)^2$ e $ C \neq I $,
então $\mu_C = (X-1)^2$: uma raiz dupla, então não \omterm{def:b2:reduction:diag}{diagonalizável}
(\cref{cor:b2:reduction:minpolycrit}).
\end{solution}

\begin{solution}{exo:b2:reduction:3}
$ X^2 - 5X + 6 = (X-2)(X-3)$: dividido com raízes simples, então $ u $ é
\omterm{def:b2:reduction:diag}{diagonalizável} (\cref{cor:b2:reduction:minpolycrit}), com
$\operatorname{Sp}(u) \subseteq \{2, 3\}$. \omterm{def:b2:reduction:eigen}{Espectros} possíveis:
$\{2\}$ ($ u = 2\,\mathrm{id}$), $\{3\}$ ($ u = 3\,\mathrm{id}$), ou
$\{2, 3\}$.

Poderes: procurar $ u^k = a_k\,\mathrm{id} + b_k\,u $. No autoespaços,
isso diz $2^k = a_k + 2b_k $ e $3^k = a_k + 3b_k $: resolver, $ b_k
= 3^k - 2^k $, $ a_k = 3\cdot2^k - 2\cdot 3^k $:
\[
u^k = (3\cdot 2^k - 2\cdot 3^k)\,\mathrm{id} + (3^k - 2^k)\, u .
\]
(Válido para todos os três \omterm{def:b2:reduction:eigen}{espectros}: as identidades são válidas em termos de \omterm{def:b2:reduction:eigen}{autovalor}.)
\end{solution}

\begin{solution}{exo:b2:reduction:4}
$ u $ \omterm{def:b2:reduction:diag}{diagonalizável}: $ P = \prod_{\lambda}(X - \lambda)$ sobre o
\omterm{def:b2:reduction:eigen}{espectro} aniquila $ u $, divisões, raízes simples. Então $ P(u|_F) =
P(u)|_F = 0$: a restrição é aniquilada por um polinômio dividido
com raízes simples, portanto \omterm{def:b2:reduction:diag}{diagonalizável}
(\cref{cor:b2:reduction:minpolycrit}).
\end{solution}

\begin{solution}{exo:b2:reduction:5}
$\chi_A = X^2 - X - 1$, raízes $\varphi $ e $\psi $ (distinto):
\omterm{def:b2:reduction:diag}{diagonalizável}, com vetores próprios $(\varphi, 1)$ e $(\psi, 1)$.
A recorrência dá $\begin{pmatrix} F_{n+1}\\ F_n \end{pmatrix} =
A^n \begin{pmatrix}1\\ 0\end{pmatrix}$. Decompor $(1, 0)$ no
vetores próprios: $(1,0) = \frac{1}{\varphi - \psi}\bigl((\varphi, 1) -
(\psi, 1)\bigr)$ com $\varphi - \psi = \sqrt5$. Candidatura $ A^n $
multiplica cada autocomponente por seu autovalorde $ n $-ésimo poder;
lendo a segunda coordenada:
\[
F_n = \frac{\varphi^n - \psi^n}{\sqrt 5} .
\]
(Verificar: $ n = 1$ dá $\frac{\varphi - \psi}{\sqrt5} = 1$.)
\end{solution}

\begin{solution}{exo:b2:reduction:6}
Deixar $ P = \prod_i (X - \mu_i)$ aniquilar $ u^2$, dividir com simples
raízes $\mu_i $ (o \omterm{def:b2:reduction:eigen}{espectro} de $ u^2$). Desde $ u $ é invertível, $0$
não é um \omterm{def:b2:reduction:eigen}{autovalor} de $ u^2$ ($\det u^2 = (\det u)^2 \neq 0$), então
tudo $\mu_i \neq 0$. Então
\[
Q(X) = \prod_i (X^2 - \mu_i) = \prod_i (X - \sqrt{\mu_i})(X +
\sqrt{\mu_i})
\]
aniquila $ u $: $\;Q(u) = \prod_i (u^2 - \mu_i\,\mathrm{id}) =
P(u^2) = 0$. Suas raízes $\pm
\sqrt{\mu_i}$ (raízes quadradas complexas) são distintas aos pares porque
o $\mu_i $ são distintos e diferentes de zero ($\sqrt{\mu_i} = -\sqrt{\mu_j}$
daria $\mu_i = \mu_j $). Divisão + raízes simples: $ u $ é
\omterm{def:b2:reduction:diag}{diagonalizável}.

Contra-exemplo sem invertibilidade: $ u = \begin{pmatrix} 0 & 1\\
0 & 0\end{pmatrix}$: $ u^2 = 0$ é \omterm{def:b2:reduction:diag}{diagonalizável}, $ u $ não é.
\end{solution}

\begin{solution}{exo:b2:reduction:7}
$ A = 2I + N $ com $ N $ a mudança ($ N e_2 = e_1$, $ Ne_3 = e_2$),
$ N^3 = 0$, $ N^2 = E_{13}$: esse \emph{é} o \omterm{thm:b2:reduction:dunford}{Decomposição de Dunford}
($2I $ diagonal, $ N $ nilpotentes, eles comutam; a singularidade faz com que
aquele). Binômio com termos comutados:
\[
A^k = 2^k I + k 2^{k-1} N + \binom k2 2^{k-2} N^2
= \begin{pmatrix}
2^k & k2^{k-1} & \binom k2 2^{k-2}\\
0 & 2^k & k2^{k-1}\\
0 & 0 & 2^k
\end{pmatrix},
\]
\[
\eu^{tA} = \eu^{2t}\Bigl(I + tN + \frac{t^2}{2}N^2\Bigr)
= \eu^{2t}\begin{pmatrix}
1 & t & t^2/2\\
0 & 1 & t\\
0 & 0 & 1
\end{pmatrix}.
\]
\end{solution}

\begin{solution}{exo:b2:reduction:8}
$ X^k - 1$ aniquila $ A $ e se divide $\C $ com o $ k $ distinto
raízes $\eu^{2\iu\pi j/k}$: $ A $ é \omterm{def:b2:reduction:diag}{diagonalizável}
(\cref{cor:b2:reduction:minpolycrit}) e seu valores próprios, raízes de
$ X^k - 1$, são $ k $-as raízes da unidade.

Se além disso $ A $ é semelhante a uma matriz triangular com unidade
diagonal: todos valores próprios igual $1$, e $ A $, \omterm{def:b2:reduction:diag}{diagonalizável} com
único \omterm{def:b2:reduction:eigen}{autovalor} $1$, é $ P\,I\,P^{-1} = I $.
\end{solution}

\begin{solution}{exo:b2:reduction:9}
Escrever $ E = \bigoplus_\lambda E_\lambda(u)$
(\cref{thm:b2:reduction:diagcrit}). Cada $ E_\lambda(u)$ é
$ v $-estável: para $ x \in E_\lambda $, $ u(v(x)) = v(u(x)) = \lambda
v(x)$. A restrição de $ v $ a $ E_\lambda(u)$ é \omterm{def:b2:reduction:diag}{diagonalizável}
(\cref{exo:b2:reduction:4}): escolha uma base de $ E_\lambda(u)$ feito
de $ v $-vetores próprios. Concatenando essas bases acima de tudo $\lambda $
dá uma base de $ E $ cujos vetores são vetores próprios de \emph{ambos}
$ u $ (por adesão em $ E_\lambda(u)$) e $ v $ (por construção).
\end{solution}

\begin{solution}{exo:b2:reduction:10}
($\Rightarrow $) Deixar $ u $ ser \omterm{def:b2:reduction:diag}{diagonalizável} e $ F $ estável. Então
$ u|_F $ é \omterm{def:b2:reduction:diag}{diagonalizável} (\cref{exo:b2:reduction:4}): $ F $ tem um
base de vetores próprios, que se estende, dentro de cada global autoespaço
$ E_\lambda $, numa base de $ E_\lambda $ (teorema da base incompleta
dentro $ E_\lambda $, começando pela parte $ F $ a base está mentindo
lá --- nota $ F = \bigoplus_\lambda (F \cap E_\lambda)$ desde
$ u|_F $ é \omterm{def:b2:reduction:diag}{diagonalizável}). Os vetores adicionados abrangem um estável
suplemento (cada um está em algum $ E_\lambda $, então seu intervalo é
$ u $-estável).

($\Leftarrow $) Deixar $ F = \sum_\lambda E_\lambda(u)$ (um estábulo
subespaço) e $ G $ um suplemento estável. Se $ G \neq \{0\}$: $\chi_{u|_G}$
se divide $\C $, então $ u|_G $ tem um vetor próprio $ x \in G $
(\cref{thm:b2:reduction:trigonalization} ou diretamente a existência
de uma raiz); mas cada vetor próprio de $ u $ reside em $ F $, então $ x \in F
\cap G = \{0\}$: contradição. Por isso $ G = \{0\}$ e $ E = F $: o
autoespaços preencher $ E $, ou seja\ $ u $ é \omterm{def:b2:reduction:diag}{diagonalizável}.
\end{solution}

\begin{solution}{exo:b2:reduction:11}
Trigonalizar: $ A = PTP^{-1}$, $ T = \begin{pmatrix} \lambda & c\\ 0 &
\mu\end{pmatrix}$, $\abs\lambda, \abs\mu < 1$. Então $ A^k =
PT^kP^{-1}$, e basta que $ T^k \to 0$.

\emph{Distinto valores próprios:} indução dá
\[
T^k = \begin{pmatrix}
\lambda^k & c\,\dfrac{\lambda^k - \mu^k}{\lambda - \mu}\\[4pt]
0 & \mu^k
\end{pmatrix},
\]
e cada entrada tende a $0$ ($\abs{\lambda}^k, \abs\mu^k \to 0$).

\emph{Igual valores próprios ($\mu = \lambda $):} $ T = \lambda I + cE_{12}$
e $ T^k = \lambda^k I + k\lambda^{k-1}cE_{12}$; a entrada
$ k\lambda^{k-1} \to 0$ desde $\abs\lambda < 1$ (batidas geométricas
polinômio). Em ambos os casos $ T^k \to 0$ entrada, portanto $ A^k =
PT^kP^{-1} \to 0$ (multiplicação de matrizes por fixo $ P, P^{-1}$ é
contínuo nas entradas --- cada entrada do produto é fixa
combinação linear).
\end{solution}

\begin{solution}{exo:b2:reduction:12}
$\ker u $ tem dimensão $ n - 1$ (classificação - nulidade), então $0$ é um
\omterm{def:b2:reduction:eigen}{autovalor} de multiplicidade geométrica $ n - 1$, e $\chi_u $ é
divisível por $ X^{n-1}$
(\cref{def:b2:reduction:charpoly}: geométrico $\leq $ algébrico).
Escrever $\chi_u = X^{n-1}(X - \alpha)$; o coeficiente de
$ X^{n-1}$ ser $-\operatorname{tr} u $, obtemos $\alpha =
\operatorname{tr} u $: $\chi_u = X^{n-1}(X - \operatorname{tr}
u)$.

Se $\operatorname{tr} u \neq 0$: o \omterm{def:b2:reduction:eigen}{autovalor}
$\operatorname{tr} u $ é uma raiz de $\chi_u $, então carrega um
vetor próprio; o autoespaços para $0$ e $\operatorname{tr} u $
tem dimensões $ n - 1$ e $\geq 1$, somando $\geq n $: eles
preencher $ E $, e $ u $ é \omterm{def:b2:reduction:diag}{diagonalizável}
(\cref{thm:b2:reduction:diagcrit}). Se $\operatorname{tr} u = 0$:
por \cref{exo:b2:linalg:5}, $ u^2 = (\operatorname{tr} u)u = 0$
com $ u \neq 0$: $ u $ é um nilpotente diferente de zero e um \omterm{def:b2:reduction:diag}{diagonalizável}
nilpotente é zero (\cref{prop:b2:reduction:nilpotent}): não
\omterm{def:b2:reduction:diag}{diagonalizável}.
\end{solution}

\begin{solution}{pb:b2:reduction:1}
\textbf{1.} O primeiro $ k - 1$ coordenadas de $ Cv_n $ são $ u_{n+1},
\dots, u_{n+k-1}$ (os deslocamentos superdiagonais), e o último é
$ a_0 u_n + \dots + a_{k-1}u_{n+k-1}$. Então $ v_{n+1} = Cv_n $ detém
para todos $ n $ se as últimas coordenadas corresponderem a todos $ n $, ou seja,\ se
$(\mathcal R)$ segura. Iterando, $ v_n = C^nv_0$.

\textbf{2.} Expandir $ D_k(X) = \det(XI_k - C)$ ao longo do primeiro
coluna: as duas entradas diferentes de zero são $ X $ (posição $(1,1)$) e
$-a_0$ (posição $(k,1)$). O primeiro menor é $ D_{k-1}$ em forma
para os coeficientes $ a_1, \dots, a_{k-1}$; o segundo menor é
triangular superior com diagonal $-1$: determinante $(-1)^{k-1}$,
com sinal $(-1)^{k+1}$ da posição. Indução ativada $ k $
(base $ k = 1$: $ X - a_0$) dá
\[
D_k(X) = X\bigl(X^{k-1} - a_{k-1}X^{k-2} - \dots - a_1\bigr) -
a_0 = P(X).
\]
Para $\mu_C $: desde $ Q(C^{\mathsf T}) = Q(C)^{\mathsf T}$ para qualquer
polinômio, $ C $ e $ C^{\mathsf T}$ tem o mesmo \omterm{def:b2:linalg:annihilator}{aniquiladores},
daí o mesmo polinômio mínimo. Para $ C^{\mathsf T}$: o
colunas lidas $ C^{\mathsf T}e_1 = e_2$, \dots, $ C^{\mathsf
T}e_{k-1} = e_k $, então $(e_1, C^{\mathsf T}e_1, \dots, (C^{\mathsf
T})^{k-1}e_1)$ é a base canônica: grátis. Um polinômio $ Q
\neq 0$ de grau $< k $ então tem $ Q(C^{\mathsf T})e_1 \neq 0$
(é uma combinação não trivial de vetores de base): $\deg\mu \geq
k $. Como $\mu \mid \chi = P $ com $\deg P = k $: $\mu_C = P $.

\textbf{3.} Para $ v = (1, \lambda, \dots, \lambda^{k-1})^{\mathsf
T}$: linhas $1$ para $ k-1$ de $ Cv $ dar $\lambda, \lambda^2, \dots,
\lambda^{k-1}$, ou seja\ $\lambda $ vezes o primeiro $ k - 1$ entradas
de $ v $; a última linha dá $\sum_m a_m\lambda^m = \lambda^k -
P(\lambda) = \lambda^k = \lambda\cdot\lambda^{k-1}$. Então $ Cv =
\lambda v $. Por outro lado, as equações $(Cx)_i = \lambda x_i $ para
$ i < k $ ler $ x_{i+1} = \lambda x_i $: qualquer vetor próprio é
proporcional a $ v $ --- todo autoespaço tem dimensão exatamente
$1$. \omterm{def:b2:reduction:diag}{Diagonalizável} se o autoespaço dimensões somam a $ k $
(\cref{thm:b2:reduction:diagcrit}) se houver $ k $ distinto
valores próprios se $ P $ tem $ k $ raízes distintas (a valores próprios são
as raízes de $\chi_C = P $).

\textbf{4.} Cada $(\lambda_i^n)_n $ resolve $(\mathcal R)$:
$\lambda_i^{n+k} = \lambda_i^n\,\lambda_i^k =
\lambda_i^n\sum_m a_m\lambda_i^m $. Liberdade: um desaparecimento
combinação $\sum_i c_i\lambda_i^n = 0$ para $ n = 0, \dots, k-1$
é um sistema Vandermonde (\cref{exo:b2:linalg:11}) no $ c_i $:
todos $ c_i = 0$. O espaço de solução tem dimensão $ k $ (pergunta 6,
cuja prova é elementar e independente): $ k $ soluções gratuitas
formam uma base e as coordenadas são únicas.

\textbf{5.} $ P = X^2 - X - 6 = (X - 3)(X + 2)$: solução geral
$ u_n = A\,3^n + B(-2)^n $. Condições iniciais: $ A + B = 1$, $3A -
2B = 8$: $ A = 2$, $ B = -1$:
\[
u_n = 2\cdot 3^n - (-2)^n .
\]
(Verificar: $ u_2 = u_1 + 6u_0 = 14$ e $2\cdot9 - 4 = 14$.)

\textbf{6.} $ P(S)\bigl((u_n)\bigr)$ é a sequência $ n \mapsto
u_{n+k} - a_{k-1}u_{n+k-1} - \dots - a_0u_n $: desaparece se
$(\mathcal R)$ vale, então o conjunto solução é $\ker P(S)$, um
subespaço. O aplicação $\ker P(S) \to \C^k $, $ u \mapsto (u_0, \dots,
u_{k-1})$, é linear, injetivo (a recorrência determina
$ u_{k}, u_{k+1}, \dots $ desde o primeiro $ k $ valores, por indução)
e sobrejetivo (definir $ u_n $ recursivamente a partir de quaisquer dados iniciais):
dimensão $ k $.

\textbf{7.} A prova de \cref{thm:b2:reduction:kernels} usa
apenas: a identidade Bézout em $\C[X]$, e o fato de que
polinômios em um deslocamento de endomorfismo fixo. Nenhum dos dois menciona
a dimensão do espaço ambiente: o lema é válido literalmente para
$ S \in \mathcal{L}(\mathcal{S})$. Por isso
\[
\ker P(S) = \bigoplus_{i=1}^{r}
\ker\,(S - \lambda_i\,\mathrm{id})^{m_i}.
\]

\textbf{8.} Para $ Q \in \C[X]$: $(S -
\lambda)\bigl(Q(n)\lambda^n\bigr)_n $ tem $ n $-ésimo termo
$ Q(n{+}1)\lambda^{n+1} - \lambda Q(n)\lambda^n =
\lambda^{n+1}(\Delta Q)(n)$, com $\Delta Q = Q(X{+}1) - Q(X)$
de grau $\deg Q - 1$ (os termos iniciais são cancelados). Iterando,
$(S - \lambda)^m\bigl(Q(n)\lambda^n\bigr) =
\bigl(\lambda^{n+m}(\Delta^m Q)(n)\bigr)_n $, e $\Delta^m Q = 0$
quando $\deg Q \leq m - 1$: o conjunto da direita está contido no
núcleo. É um subespaço de dimensão $ m $: as sequências
$(n^j\lambda^n)_n $, $0 \leq j < m $, são gratuitos, pois
$\sum_j c_j n^j\lambda^n = 0$ para todos $ n $ forças (dividindo por
$\lambda^n \neq 0$) o polinômio $\sum_j c_jX^j $ desaparecer em
cada $ n \in \N $, portanto, ser zero. Por outro lado $\dim\ker(S -
\lambda)^m \leq m $: expandindo $(S - \lambda)^m = \sum_j
\binom mj(-\lambda)^{m-j}S^j $, a equação $(S - \lambda)^m u =
0$ é uma recorrência linear de ordem $ m $ (coeficiente líder
$1$), então $ u $ é determinado por $ u_0, \dots, u_{m-1}$ como em
questão 6. Conclui-se a igualdade de dimensões.

\textbf{9.} Combine as questões 7 e 8: toda solução se decompõe
exclusivamente como uma soma de elementos do $\ker(S -
\lambda_i)^{m_i}$, ou seja\ $ u_n = \sum_i Q_i(n)\lambda_i^n $ com
$\deg Q_i \leq m_i - 1$; o $ Q_i $ são únicos porque
a decomposição é direta e, dentro de cada soma, o
coeficientes de $ Q_i $ são coordenadas na base
$(n^j\lambda_i^n)_j $ (questão 8). Verificação de sanidade nas dimensões:
$\sum_i m_i = k $.

\textbf{10.} $ P = X^2 - 4X + 4 = (X - 2)^2$: soluções $(a +
bn)2^n $. Dados iniciais: $ a = 1$, $2(a + b) = 0$, então $ b = -1$:
\[
u_n = (1 - n)\,2^n .
\]
Verificar: $ u_2 = 4u_1 - 4u_0 = -4$, e $(1 - 2)\cdot4 = -4$.

\textbf{11.} Escrever $ u_n = \lambda_1^n\bigl(c_1 + \sum_{i\geq2}
c_i(\lambda_i/\lambda_1)^n\bigr)$; cada razão tem módulo $< 1$,
então o colchete tende a $ c_1 \neq 0$: $ u_n \sim c_1\lambda_1^n $.
Em particular $ u_n \neq 0$ para grande $ n $, e
\[
\frac{u_{n+1}}{u_n} =
\lambda_1\,\frac{c_1 + o(1)}{c_1 + o(1)} \longrightarrow
\lambda_1 .
\]

\textbf{12.} Calcular:
\[
q(a_{n+1}, b_{n+1}) = (a_n + 2b_n)^2 - 2(a_n + b_n)^2
= -a_n^2 + 2b_n^2 = -q(a_n, b_n).
\]
Com $ q(a_0, b_0) = 1 - 2 = -1$: $ a_n^2 - 2b_n^2 = (-1)^{n+1}$.
Estruturalmente: $ q(a, b) = (a - \sqrt2\,b)(a + \sqrt2\,b)$ e
o aplicação linear $ M $ multiplica o fator $ a + \sqrt2 b $ por $1 +
\sqrt2$ e o fator $ a - \sqrt2 b $ por $1 - \sqrt2$ (calcular:
$ a_{n+1} + \sqrt2 b_{n+1} = (1 + \sqrt2)(a_n + \sqrt2 b_n)$);
o produto é multiplicado por $(1 + \sqrt2)(1 - \sqrt2) = -1 =
\det M $ cada passo.

\textbf{13.} Desde $ a_n^2 - 2b_n^2 = (a_n - \sqrt2 b_n)(a_n +
\sqrt2 b_n) = (-1)^{n+1}$,
\[
\Bigl|\frac{a_n}{b_n} - \sqrt2\Bigr|
= \frac{\abs{a_n^2 - 2b_n^2}}{b_n(a_n + \sqrt2 b_n)}
= \frac{1}{b_n(a_n + \sqrt2 b_n)} \leq \frac1{2b_n^2},
\]
usando $ a_n \geq b_n \geq 1$ (indução: ambos aumentam) então $ a_n +
\sqrt2 b_n \geq (1 + \sqrt2)b_n \geq 2b_n $. \omterm{def:b2:reduction:eigen}{Autovalores} de $ M $:
$\chi_M = X^2 - 2X - 1$, raízes $1 \pm \sqrt2$; desde $(a_0, b_0)$
tem um componente diferente de zero no dominante vetor próprio (todas as entradas
positivo), $ b_n \sim c(1 + \sqrt2)^n $ com $ c > 0$ (pergunta
11). Portanto o erro é $\asymp (1 + \sqrt2)^{-2n} = (3 +
2\sqrt2)^{-n}$: decaimento geométrico com proporção $1/(3 + 2\sqrt2) = 3
- 2\sqrt2 \approx 0.172$.

\textbf{14.} (a) Da pergunta 9: $\abs{u_n} \leq \sum_i
\abs{Q_i(n)}\abs{\lambda_i}^n \leq \bigl(\sum_i
\abs{Q_i(n)}\bigr)\rho^n $, e cada $\abs{Q_i(n)} \leq C_i
n^{m_i - 1} \leq C_i n^{m-1}$ para $ n \geq 1$: soma as constantes.
(b) Seja $\rho' = \max_{i \geq 2}\abs{\lambda_i} <
\abs{\lambda_1}$ e $ d = \deg Q_1$, coeficiente líder $ c
\neq 0$. Então $ u_n = Q_1(n)\lambda_1^n + R_n $ com $\abs{R_n}
\leq Cn^{m-1}\rho'^n $, e
\[
\frac{R_n}{Q_1(n)\lambda_1^n} = O\Bigl(n^{m-1-d}
\bigl(\rho'/\abs{\lambda_1}\bigr)^n\Bigr) \longrightarrow 0
\]
(geométrica bate polinômio). Então
\[
u_n \sim Q_1(n)\,\lambda_1^n \sim c\,n^d\lambda_1^n,
\qquad
\frac{u_{n+1}}{u_n} \longrightarrow \lambda_1
\quad\text{(since } Q_1(n{+}1)/Q_1(n) \to 1\text{)}.
\]
Verifique na questão 10: para $ u_n = (1-n)2^n $ a proporção é
\[
\frac{(-n)2^{n+1}}{(1-n)2^n} = 2\,\frac{-n}{1-n}
\longrightarrow 2 = \lambda_1 .
\]

\textbf{15.} Indução ativada $ n $. Para $ n = 1$, $ A_{ij}$ conta
caminhadas de comprimento $1$. Passo: uma caminhada longa $ n + 1$ de $ i $ a $ j $ é uma caminhada de comprimento $ n $ de $ i $ a algum vértice $\ell $
seguido por uma borda $\ell j $:
\[
\#\{\text{walks}\} = \sum_{\ell} (A^n)_{i\ell}A_{\ell j} =
(A^{n+1})_{ij}.
\]

\textbf{16.} $ J = 3\Pi $ onde $\Pi = J/3$ é a projeção em
$\operatorname{Vect}(1,1,1)$ ao longo do plano $ x + y + z = 0$
($\Pi^2 = \Pi $ desde $ J^2 = 3J $). Então $ A = J - I = 2\Pi - (I -
\Pi)$, e desde $\Pi $ e $ I - \Pi $ são complementares
projeções,
\[
A^n = 2^n\,\Pi + (-1)^n (I - \Pi),
\qquad\text{i.e.}\qquad
(A^n)_{ij} = \frac{2^n}3 + (-1)^n\Bigl(\delta_{ij} -
\frac13\Bigr),
\]
que fornece as duas fórmulas exibidas. No $ n = 2$: diagonal
$(4 + 2)/3 = 2$ (caminha $ i \to \ell \to i $ para os dois vizinhos
$\ell $); fora da diagonal $(4 - 1)/3 = 1$ (a única caminhada $ i \to
\ell \to j $ através do terceiro vértice).

\textbf{17.} Deixar $ w_n^{(0)}, w_n^{(1)}$ contar palavras admissíveis de
comprimento $ n $ terminando em $0$, resp.\ $1$. Anexando uma carta: a $0$
pode seguir qualquer coisa, um $1$ apenas um $0$:
\[
\begin{pmatrix} w_{n+1}^{(0)}\\ w_{n+1}^{(1)}\end{pmatrix}
= \begin{pmatrix} 1 & 1\\ 1 & 0\end{pmatrix}
\begin{pmatrix} w_n^{(0)}\\ w_n^{(1)}\end{pmatrix}.
\]
Resumindo, $ w_{n+2} = w_{n+1} + w_n $ (ou: condição no primeiro
carta). Com $ w_1 = 2$, $ w_2 = 3$: $ w_n = F_{n+2}$ por indução
($ F_3 = 2$, $ F_4 = 3$, mesma recorrência). Crescimento: as raízes do
$ X^2 - X - 1$ são $\varphi > \abs\psi $
(\cref{exo:b2:reduction:5}) e o $\varphi $-componente é
diferente de zero (o $ w_n $ são positivos e $\psi^n \to 0$), então questione
11 dá $ w_{n+1}/w_n \to \varphi = \frac{1 + \sqrt5}2$.

\textbf{18.} $ A = \begin{psmallmatrix} 0&1&0\\ 1&0&1\\
0&1&0\end{psmallmatrix}$. Verificar:
\[
A(1, \pm\sqrt2, 1)^{\mathsf T}
= (\pm\sqrt2, 2, \pm\sqrt2)^{\mathsf T}
= \pm\sqrt2\,(1, \pm\sqrt2, 1)^{\mathsf T},
\qquad
A(1, 0, -1)^{\mathsf T} = 0 :
\]
valores próprios $\sqrt2, -\sqrt2, 0$ ($= 2\cos\frac\pi4,
2\cos\frac{3\pi}4, 2\cos\frac\pi2$). Decompor $ e_1$ no
autobase e leia a terceira coordenada, ou use simetria: com
$ v_\pm = (1, \pm\sqrt2, 1)$, $ v_0 = (1, 0, -1)$, verifica-se $ e_1
= \frac14 v_+ + \frac14 v_- + \frac12 v_0$, então para $ n \geq 1$
\[
(A^n)_{13} = \Bigl(\tfrac14(\sqrt2)^n v_+ +
\tfrac14(-\sqrt2)^n v_- + 0\Bigr)_{\!3}
= \frac{(\sqrt2)^n + (-\sqrt2)^n}{4},
\]
zero para ímpar $ n $ (gráfico bipartido: as extremidades estão em distâncias iguais),
e $2\cdot 2^{n/2}/4 = 2^{n/2 - 1}$ até mesmo $ n $. No $ n = 4$:
$2^{1} = 2$, combinando os dois passeios $1\,2\,1\,2\,3$ e
$1\,2\,3\,2\,3$.

\textbf{19.} Caminhadas fechadas de comprimento $ n $ de $ i $ são
$(A^n)_{ii}$; resumindo $ i $ dá $\operatorname{tr}(A^n)$.
Trigonalizando $ A $ (sobre $\C $), $ A^n $ é triangular com diagonal
$\lambda_i^n $: $\operatorname{tr}(A^n) = \sum_i\lambda_i^n $.
Triângulo: $\operatorname{tr}(A^n) = 3\,\frac{2^n + 2(-1)^n}3 =
2^n + 2(-1)^n = 2^n + (-1)^n + (-1)^n $: o \omterm{def:b2:reduction:eigen}{espectro} $\{2, -1,
-1\}$, consistente com a questão 16.

\textbf{20.} As colunas de $ W^{\mathsf T}$: $ W^{\mathsf T}e_i =
e_{i-1}$ para $ i \geq 1$ e $ W^{\mathsf T}e_0 = e_{k-1}$;
reetiquetando na ordem $ e_0, e_1, \dots $ este é exatamente o
matriz companheira de $ X^k - 1$ ($ a_0 = 1$, outro $ a_m = 0$).
Pergunta 2: $\chi_{W} = \chi_{W^{\mathsf T}} = X^k - 1 =
\mu_{W}$. As raízes $\omega^j $ ($ j = 0, \dots, k-1$) são os $ k $
distinto $ k $-ésimas raízes da unidade: $ W $ é \omterm{def:b2:reduction:diag}{diagonalizável} (pergunta
3, ou \cref{exo:b2:reduction:8}: $ W^k = I $). Autovetores: $ Wf_j
= \sum_m \omega^{-jm}e_{m+1} = \sum_{m'}\omega^{-j(m'-1)}e_{m'}
= \omega^j f_j $.

\textbf{21.} Circulantes são polinômios em $ W $ e polinômios
em uma matriz fixa comutam entre si. Cada $ f_j $ é um
vetor próprio de todos os poderes: $ W^m f_j = \omega^{jm}f_j $, então
\[
Cf_j = \sum_m c_m\omega^{jm} f_j = \widehat c(\omega^j)\,f_j :
\]
a base $(f_0, \dots, f_{k-1})$ (grátis: Vandermonde no
distinto $\omega^{-j}$, \cref{exo:b2:linalg:11}) diagonaliza
todos os circulantes de uma só vez, com o declarado valores próprios.

\textbf{22.} O determinante é o produto do valores próprios
(diagonalizar): $\det C = \prod_{j}\widehat c(\omega^j)$. Para $ k
= 3$, $ c_0 = a $, $ c_1 = b $, $ c_2 = c $ e $\omega = j =
\eu^{2\iu\pi/3}$:
\[
\det C = (a + b + c)(a + bj + cj^2)(a + bj^2 + cj^4),
\]
e $ j^4 = j $: exatamente a fatoração de
\cref{exo:b2:linalg:8}.

\textbf{23.} $ M = \frac12(W + W^{-1})$ é um circulante ($ W^{-1} =
W^{k-1}$), com valores próprios $\frac12(\omega^j + \omega^{-j}) =
\cos\frac{2\pi j}k $ na mesma base $ f_j $. Coordenadas: escrever
$ x^{(0)} = \sum_j \alpha_j f_j $. As coordenadas de $ f_j $ soma para
$\sum_m \omega^{-jm}$, que é $ k $ para $ j = 0$ e $0$
caso contrário (soma geométrica com razão $\omega^{-j} \neq 1$).
Somando as coordenadas de $ x^{(0)}$: $\sum_m x^{(0)}_m =
\alpha_0\,k $, então $\alpha_0 = \frac1k\sum_m x^{(0)}_m $, a média.

\textbf{24.} $ x^{(n)} = M^nx^{(0)} = \sum_j
\alpha_j\cos^n\bigl(\tfrac{2\pi j}k\bigr)f_j $. Por estranho $ k $,
$\abs{\cos(2\pi j/k)} < 1$ para cada $ j \neq 0$ (o ângulo é
nunca $0$ ou $\pi $), então todos os termos, exceto $ j = 0$ tendem a $0$:
$ x^{(n)} \to \alpha_0 f_0$, o vetor constante igual ao
média --- a média em um anel ímpar equaliza. Para $ k = 4$ o
valores próprios são $1, 0, -1, 0$: o $ j = 2$ prazo
$\alpha_2(-1)^nf_2$ com $ f_2 = (1, -1, 1, -1)^{\mathsf T}$
oscila para sempre. A obstrução é a média alternada:
multiplicando as coordenadas de $ x^{(0)}$ por $(-1)^m $ e
somando, o mesmo cálculo da soma geométrica dá $\sum_m
(-1)^mx^{(0)}_m = 4\alpha_2$: o processo converge se e se
$ x^{(0)}_0 - x^{(0)}_1 + x^{(0)}_2 - x^{(0)}_3 = 0$ e então
converge para a média.

\textbf{25.} A matriz complementar converte uma recorrência escalar
de ordem $ k $ em uma recorrência vetorial de primeira ordem, de modo que
fórmulas fechadas tornam-se asserções sobre $ C^n $ --- redução
base (questões 1 a 5). O lema da decomposição do núcleo é
álgebra polinomial pura (Bézout mais comutação), então ela divide
$\ker P(S)$ embora $\mathcal{S}$ é de dimensão infinita
(questões 7 a 9). Dominante valores próprios governar o crescimento porque
todas as outras contribuições são geometricamente desprezíveis após
normalização --- é também por isso que o erro Pell decai no
quadrado da raiz dominante (questões 11-14). Poderes do
a contagem da matriz de adjacência caminha porque a multiplicação da matriz soma
sobre vértices intermediários, então os \omterm{def:b2:reduction:eigen}{espectros} contam passeios fechados
(questões 15 a 19). Matrizes de comutação compartilham uma base própria, e
uma base de Fourier diagonaliza toda a álgebra circulante
de uma só vez (questões 20 a 24). Cúpulas: o fundamental
teorema das recorrências lineares (questão 9); e para não negativo
matrizes, a razão pela qual raízes dominantes como $\varphi $ ou $1 +
\sqrt2$ são automaticamente reais, positivos e simples é o
Teorema de Perron-Frobenius, provado no volume do Ano 3.
\end{solution}
